\chapter{把计划当作架构控制}
\label{ch:planning}

\section{在行动时刻之前完成计划}

\begin{theorem}[计划定理]
在高澄清度时段完成的准备工作，会比临时说服自己更可靠地降低未来的激活障碍。
\end{theorem}

\begin{discussion}
计划至少通过四条通道发挥作用：
\begin{enumerate}[nosep]
  \item 它移除了目标状态中的歧义；
  \item 它缩短了行动当下仍需现场决定的链条；
  \item 它把物质线索提前放进环境里；
  \item 它可以通过日程、提醒或其他人，把外部力量导入进来。
\end{enumerate}
因此，计划并不只是某种激励性语言。它本质上是在重写未来选择发生时的结构条件。
\end{discussion}

\section{计划栈}

为便于教学，把计划分成几个层次会更清楚：
\begin{enumerate}[nosep]
  \item \textbf{结果层：} 想要达到的终局状态是什么？
  \item \textbf{会话层：} 哪一段具体会话能推动这个结果？
  \item \textbf{入口层：} 哪个确切的第一步动作会启动这段会话？
  \item \textbf{环境层：} 什么样的物理布置能够降低进入摩擦？
  \item \textbf{触发层：} 预期中的决策窗口会在何时、何地出现？
\end{enumerate}

\begin{proposition}[二元执行优势]
当一个计划把未来行动压缩为二元选择---执行或拒绝---时，它消耗的意志能量会少于那种仍然要求现场重新设计的计划。
\end{proposition}

\section{章级工作流}

\begin{templatebox}
\textbf{最小转移计划}
\begin{enumerate}[nosep]
  \item 写出目标状态。
  \item 写出最可能出现的决策窗口。
  \item 写下第一个可见动作。
  \item 把所需物品提前放入未来环境。
  \item 定义这次会话的停止规则。
\end{enumerate}
\end{templatebox}

\section{未经规划的意图为何会失败}

\begin{corollary}[未经规划的行动者会追随当前惯性]
如果没有任何准备去降低未来障碍项，那么下一次决策中的默认赢家就会是当前状态，而不是口头上更偏好的那个状态。
\end{corollary}

\begin{example}[下班后锻炼]
“我下班后应该去锻炼”这句话，还不能算一个可用的计划。它没有解决地点、背包、动作顺序、训练内容、时长和退出条件。这些没有回答的问题，都是隐藏起来的障碍项。真正的计划，会在上游把它们逐一消解。
\end{example}

\section{推荐阅读}

\begin{readingbox}
\textbf{基础}
\begin{itemize}[nosep]
  \item Stanislas Dehaene, \emph{Consciousness and the Brain}，可用于理解结构化通达与全局协调的更大逻辑。
\end{itemize}
\textbf{现代研究}
\begin{itemize}[nosep]
  \item 关于意识状态特征与稳定性的研究，有助于说明为什么计划应被视为状态工程，而不是打气式劝说。
\end{itemize}
\textbf{前沿}
\begin{itemize}[nosep]
  \item 要把计划理解为一种架构控制，并不需要诉诸前沿物理理论。
\end{itemize}
\end{readingbox}

\begin{summarybox}
这一章把计划转化为基础设施。它对整个项目的价值，在于给后续案例和练习提供了一套稳定模板：识别窗口、降低歧义、预装环境，并把未来选择压缩下来。
\end{summarybox}
